\subsection{The reverse cyclic-four case and the exact smallest height-two index}

\begin{lemma}[The two cyclic-four wild families and their cubic fields]
\label{lem:B-RD42-normal}
\status{Proved. Complete root and independent original-source, proof and exact-certificate audits passed.}
Let $K$ be complete of mixed characteristic $(0,3)$, with algebraically
closed residue field and $v_K(3)=2$. Fix $\pi^2=3$ and $i^2=-1$ in $K$.
A potentially supersingular elliptic curve with irreducible cyclic-four
three-torsion representation and faithful special inertia of order six
is $K$-isomorphic to
\[
 M_b:\ y^2=x^3-3x+3b,\qquad b\in\OO_K^\times.
                                                        \tag{B-RD42.1}
\]
A curve with the same residual condition and faithful special inertia
of order three is the quadratic $\pi$-twist $M_b^{(\pi)}$ of such a curve.
The minimal discriminant valuations are respectively six and twelve.

Put $L_b=K(r)$, where $r^3-3r+3b=0$. This is a cyclic cubic extension,
and $M_b^{(\pi)}$ has good reduction over $L_b$, with faithful special
inertia of order three. The curve $M_b$ has good reduction over
$L_b(\sqrt\pi)$, with faithful special inertia of order six. Moreover
\[
 L_b=L_{b'}\quad\Longleftrightarrow\quad
                   \overline {b'}=\overline b\ \hbox{or}\ -\overline b.
                                                        \tag{B-RD42.2}
\]
\end{lemma}
\begin{proof}
The discriminant bound in Proposition~\ref{prop:B-RD36} and
Proposition~\ref{prop:B-RD38} give
$0\le\delta\le23$, $6\mid\delta$, and $v_K(j)=6$.
Thus $v_K(c_4)=(\delta+6)/3$. Passing to the short equation
$y^2=x^3+Ax+B$, with $A=-c_4/48$ and $B=-c_6/864$, gives an
integral equation when $\delta=6,12,18$. Indeed
$c_6^2=c_4^3-1728\Delta$ gives respectively
$v_K(c_6)\ge6,9,12$, so the pairs of coefficient valuations satisfy
\[
 (v_K(A),v_K(B))=(2,\ge0),\ (4,\ge3),\ (6,\ge6).
\]
The last equation scales down by $x=\pi^2X,y=\pi^3Y$, contradicting
minimality. The tame-order formula of Proposition~\ref{prop:B-RD36}
therefore forces $\delta=6$ for inertia order six and $\delta=12$
for order three; $\delta=0$ would mean good reduction.

In the first case write $A=3a$, $a$ a unit. The four three-torsion
lines have $x$-coordinates given by
\[
                     x^4+6ax^2+4Bx-3a^2=0.
                                                        \tag{B-RD42.3}
\]
If $v_K(B)=0$, its Newton polygon has one root of valuation two
and three of valuation zero. If $v_K(B)=1$, it has one root of
valuation one and three of valuation $1/3$. In each case the unique
root valuation gives a Galois-stable three-torsion line, contrary
to irreducibility of the cyclic-four representation. Hence
$v_K(B)\ge2$. If $v_K(B)\ge3$, the $\pi$-twist has coefficient
valuations at least $(4,6)$; scaling by $\pi$ makes its discriminant
a unit. The original curve would acquire good reduction over a
quadratic extension, contrary to its wild inertia of order three.
Thus $B=3b_0$ with $b_0$ a unit. A fourth root of $-a$ exists in $K$
by Hensel's lemma and normalizes $a$ to $-1$, giving (B-RD42.1).

For the order-three curve, twist its integral short equation by
$\pi$ and scale by $\pi$. The result has discriminant valuation six,
coefficient valuations $(2,\ge0)$, cyclic-four residual image, and
faithful special inertia of order six. (The unique nontrivial tame
quadratic character multiplies its central tame action by $-1$.)
The preceding argument applies. Twisting back proves the assertion.

For completeness the good fields can be exhibited directly.
The Newton polygon of $r^3-3r+3b$ gives $v_K(r)=2/3$, so its
root field has degree three. Its discriminant is
$27(4-9b^2)$, a square in $K$; it is therefore cyclic. All three
pairwise root differences have valuation one, since Galois acts
transitively on the three unordered pairs and the discriminant has
valuation six. The $\pi$-twist of (B-RD42.1) has equation
$y^2=x^3-9x+\pi^5b$. The substitution
$x=\pi^2X+\pi r$, $y=\pi^3Y$ gives
\[
              Y^2=X^3+\pi rX^2+(r^2-1)X.
                                                        \tag{B-RD42.4}
\]
Its reduction is $Y^2=X^3-X$, a good supersingular curve.
A nontrivial $\sigma\in\operatorname{Gal}(L_b/K)$ changes $X$ by
$(r-\sigma r)/\pi$, whose residue is nonzero, and fixes $Y$.
Thus its action on the special curve is nontrivial of order three.
Untwisting requires $\sqrt\pi$; the added quadratic action is
$[-1]$, proving both faithful inertia assertions. In particular
these displayed good fields are the minimal good fields.

One can also check the ramification directly: $r^2/\pi$ is a
uniformizer of $L_b$, and
\[
 v_{L_b}\bigl(\sigma(r^2/\pi)-r^2/\pi\bigr)=2.
\]
The cubic lower break is one. This observation is not needed for
the equality criterion below.

Choose the square root $v_b=\sqrt{1-9b^2/4}$ congruent to one, and put
\[
                 D_b=v_b+\frac{3i}{2}b.
                                                        \tag{B-RD42.5}
\]
If $u^3=iD_b$, then $u^3+u^{-3}=-3b$, so $u+u^{-1}$ is a root
of the cubic. A cube root of $i$ lies in $K$, and
\[
                    D_b\equiv1+\frac{3i}{2}b\pmod{\pi^3}.
                                                        \tag{B-RD42.6}
\]
This element is not a cube: its deviation from one has valuation
two, whereas a cube congruent to one has deviation at least three.
Consequently $K(\sqrt[3]{D_b})=L_b$. By
Lemma~\ref{lem:B-RD39-units}, every element of $1+\pi^3\OO_K$ is a
cube. The Kummer classes of $D_b$ and $D_{b'}$ therefore generate
the same line in $K^\times/(K^\times)^3$ exactly when
$\overline{b'}=\pm\overline b$. This proves (B-RD42.2).
\end{proof}

\begin{lemma}[The reverse marked cocycle identifies the cubic good fields]
\label{lem:B-RD42-reverse}
\status{Proved. Complete root and independent original-source, proof and exact-certificate audits passed.}
Suppose that $M_b[9]\simeq M_{b'}^{(\pi)}[9]$ is antisymplectic,
and that their common residual representation is irreducible cyclic
of order four. Then $L_b=L_{b'}$.
\end{lemma}
\begin{proof}
Let $L=K(M_b[3])$ and let $F$ be its quadratic subfield. The residual
isomorphism gives the same fields on the other side. Choose symplectic
three-torsion markings with a fixed Weil-pairing value $\zeta$. Since
the given isomorphism reverses the pairing, its matrix can be taken
to be $R=\operatorname{diag}(1,-1)$ between these markings. Choose
the cyclic-four generator as
$S=\left(\begin{smallmatrix}0&-1\\1&0\end{smallmatrix}\right)$
on the first side. On the second side its matrix is
$RSR^{-1}=S^{-1}=-S$. This is precisely Fisher's reverse-marking
involution $\iota=\operatorname{diag}(1,-1)$; his inverse-transpose
twist dictionary is \citet[Section~3, p.~7, and Lemma~3.2, pp.~7--8]
{Fisher2015Nine}.

The projective action on the two Hesse parameters is the same
involution $h(t)=(t+2)/(t-1)$. Write $z=t-1$, $z'=t'-1$ and
$u=z/\pi$, $u'=z'/\pi$. As in Proposition~\ref{prop:B-RD40},
\[
 N_{F/K}(z)=N_{F/K}(z')=3,\qquad
 N_{F/K}(u)=N_{F/K}(u')=1,\qquad
 \overline u,\overline{u'}\in\{1,-1\}.
                                                        \tag{B-RD42.7}
\]
The actual Fourier lift for $S$ has differential multiplier
$\kappa(t)=c/(t-1)$, where $c=\zeta-\zeta^2$. The lift for $S^{-1}$
is its central negative: its multiplier is $-\kappa(t)$.
Indeed $S^{-1}=-S$, the central element is elliptic inversion,
and the exact normalization and identity
$\kappa(t)\kappa(h(t))=-1$ were proved in
Proposition~\ref{prop:B-RD40}. Thus no change of plane-residue
normalization or parameter-dependent factor is being suppressed.

Let $\alpha,\alpha'$ compare the respective universal Hesse
invariant differentials with invariant differentials over $K$.
The analogue of (B-RD40.7), now with opposite actual lifts, is
\[
       \frac{\tau(\alpha'/\alpha)}{\alpha'/\alpha}
                 =-\frac{z'}z\quad\hbox{or its inverse},
                                                        \tag{B-RD42.8}
\]
for a generator $\tau$ of $\operatorname{Gal}(L/K)$.
The discriminant valuations are six and twelve, so the invariant
section $\Delta\omega^{12}$ gives
$v_L(\alpha')-v_L(\alpha)=2$. Tame leading coefficients therefore
make the residue of the left side equal to $-1$. It follows that
\[
                            \overline{u'}=\overline u.
                                                        \tag{B-RD42.9}
\]

Over $F$ twist both curves by the same quadratic character of $L/F$.
Their three-torsion becomes rational, the Hesse curves are $E_t,E_{t'}$,
and the reverse nine-congruence persists. Fisher's reverse equation
(11), with its essential numerator three, gives at $\theta=1$
\[
                          u'u=\frac{(t'-1)(t-1)}3
                                    \in(F^\times)^3.
                                                        \tag{B-RD42.10}
\]
Here we use exactly the full-level marking convention of
\citet[Section~2, pp.~4--5; Theorem~4.3 and its proof, pp.~9--10]
{Fisher2015Nine}. By (B-RD42.9), this cube is congruent to one.
Since $v_F(3)=4$, expansion of $(1+w)^3$ shows
$v_F(u'u-1)\ge3$. Set
\[
 T=z+3/z=\pi(u+u^{-1}),\qquad
 T'=z'+3/z'=\pi(u'+u'^{-1})\in K.
\]
The exact identity
\[
       T'-T=\pi(u'-u^{-1})\left(1-\frac u{u'}\right)
                                                        \tag{B-RD42.11}
\]
and (B-RD42.9) give $v_F(T'-T)\ge2+3+1=6$, hence
\[
                              v_K(T'-T)\ge3.
                                                        \tag{B-RD42.12}
\]

We recover the residues of the normal-form parameters from these
traces. Eliminating $z^2-Tz+3=0$ from the Hesse $j$-formula gives
\[
 j(E_t)=27\left(\frac{T(T+4)}{T+3}\right)^3,
 \qquad
 j(M_b)=\frac{6912}{4-9b^2}.
                                                        \tag{B-RD42.13}
\]
These are exact rational identities, and twisting does not change
$j$. Put $H=T(T+4)/(T+3)$. If $\overline u=\epsilon\in\{1,-1\}$,
then the identity
\[
 u+u^{-1}-2\epsilon=\frac{(u-\epsilon)^2}{u}
\]
and $v_F(\pi)=2$ give $v_F(T-2\epsilon\pi)\ge4$. Since
$T\in K$, we may write $T=2\epsilon\pi+\pi^2a$ with
$a\in\OO_K$. Also
\[
 H-4=\frac{T^2-12}{T+3}\in\pi^2\OO_K,\qquad
 H^3=\frac{64}{1-9b^2/4}.
\]
Writing $H=4+\pi^2d$ and comparing the latter identity modulo
$\pi^5$ gives $d\equiv b^2\pmod\pi$. The former identity gives
$d\equiv2a\pmod\pi$. Consequently
\[
                T\equiv2\epsilon\pi+\frac{\pi^2}{2}b^2
                                      \pmod{\pi^3}.
                                                        \tag{B-RD42.14}
\]
The same formula holds for $T'$ and $b'$ with the same $\epsilon$
by (B-RD42.9). Equation (B-RD42.12) therefore gives
$\overline{b'}^{\,2}=\overline b^{\,2}$. In characteristic three
this says $\overline{b'}=\pm\overline b$, and
Lemma~\ref{lem:B-RD42-normal} proves $L_b=L_{b'}$.
\end{proof}

\begin{theorem}[B-RD42: the exact least height-two index at $p=3,e=2$]
\label{thm:B-RD42}
\status{Proved. Complete root and independent original-source, proof and exact-certificate audits passed.}
For the full class of height-two groups in B.4 satisfying the manuscript
monodromy hypothesis,
\[
                              \boxed{\nu_{\min}(3,2,2)=2}.
                                                        \tag{B-RD42.15}
\]
Equivalently, $27$-torsion determines the requested character on
prime-to-three inertia elements, while nine-torsion does not always do so.
This is a sharp small-height case; no general least-index formula is
asserted for other heights or ramification indices.
\end{theorem}
\begin{proof}
The lower bound is the explicit nine-torsion example of
Proposition~\ref{prop:B-RD29-p3-level-two}. The reduction and bounds in
Theorem~\ref{thm:B-RD34} and Proposition~\ref{prop:B-RD35} reduce
any putative failure at $27$ to two potentially supersingular elliptic
curves after completed maximal unramified base change. These steps
already cover arbitrary perfect residue fields and all the allowed
height-two monodromy filtrations. The completed maximal unramified
base change preserves inertia and absolute ramification; its algebraically
closed residue field makes every unit a square, so $\pi^2=3$ and
$\mu_3\subset K$ can be arranged without a further ramified extension.

Propositions~\ref{prop:B-RD37} and \ref{prop:B-RD41} exclude the
scalar and unipotent residual cases. By
Propositions~\ref{prop:B-RD38} and \ref{prop:B-RD40}, the sole remaining
case is an antisymplectic comparison with cyclic-four residual image,
individual faithful special-inertia orders three and six, and joint
faithful inertia order eighteen. Lemma~\ref{lem:B-RD42-normal}
writes this pair, after interchanging it, as $M_b,M_{b'}^{(\pi)}$.
Restriction of the normalized antisymplectic $27$-isomorphism to
nine-torsion satisfies Lemma~\ref{lem:B-RD42-reverse}; hence
$L_b=L_{b'}$. Both curves then have good reduction over the same
extension $L_b(\sqrt\pi)$ of degree six. Their joint faithful special
inertia has order at most six, contradicting the required order
eighteen. This excludes the last case and proves (B-RD42.15).
\end{proof}

\paragraph{Exact certificate and scope.}
The program \path{code/verify_c4_final_exclusion.py} checks the universal
normal-form and good-model invariants, the Cardano and Kummer identities,
the trace elimination and truncated expansions, the opposite marked
matrices, and the integer valuation cases used in the proof. Its output
is retained separately, with author runtime $0.1415409550027107$ seconds.
The initial structural-equality check and traceback are retained in
\path{failed_candidates/B42_certificate_structural_equality_v1.py} and
its adjacent output; the corrected certificate tests exact polynomial
differences. The preceding proof uses the algebraic-closure
residue hypothesis only in the established reduction and the explicitly
stated Hensel and unit-cube steps. The finite-elimination procedure of
B-RD24 remains unevaluated; it is not used to obtain this exact index.

\paragraph{Independent audits.}
The complete proof and exact certificate were independently checked by
two reviewers. Their author-output-preserving reruns passed in
$0.11300166400178568$ and $0.10002119700220646$ seconds, respectively.
The reverse-marking convention was checked against Fisher's original
Section~3 and the full proof of Theorem~4.3. This establishes only the
explicit small instance (B-RD42.15); the general least-index problem
in B.4 remains open.
